\section{Space Complexity --- Logspace --- ISCount Math}
%%% generated by the blueprint pipeline from TCSlib.Complexity.SpaceComplexity.Logspace.ISCountMath
\subsection*{Space Complexity --- Logspace --- ISCount Math: Introduction}

Declarations of the Lean module \texttt{TCSlib.Complexity.SpaceComplexity.Logspace.ISCountMath}.

\subsection*{Space Complexity --- Logspace --- ISCount Math: Results}

\begin{lemma}[Binary size under halving]
\lean{Complexity.LogProg.ISCount.size_half}
\label{Complexity.LogProg.ISCount.size_half}
\leanok
\difficulty{4}
For every nonzero natural number $m$, the binary size of $m$ is one greater than the
binary size of the integer quotient $m/2$.
\end{lemma}

\begin{lemma}[Logarithmic space bound]
\lean{Complexity.LogProg.ISCount.logSpace_le_size_succ}
\label{Complexity.LogProg.ISCount.logSpace_le_size_succ}
\leanok
\uses{Complexity.logSpace}
\difficulty{4}
For every natural number $n$, define its logarithmic space measure by $\log_2 n+1$,
where $\log_2 n$ is the natural-number logarithm. This measure is at most the binary
size of $n$ plus one.
\end{lemma}

\begin{lemma}[Binary size power bound]
\lean{Complexity.LogProg.ISCount.lt_two_pow_size_succ}
\label{Complexity.LogProg.ISCount.lt_two_pow_size_succ}
\leanok
\difficulty{2}
For every natural number $n$, one has $n<2^{\operatorname{size}(n)+1}$, where
$\operatorname{size}(n)$ is its binary size.
\end{lemma}

\begin{definition}[Membership by an initial walk]
\lean{Complexity.LogProg.ISCount.Mem}
\label{Complexity.LogProg.ISCount.Mem}
\leanok
\uses{Complexity.WalkLe}
Let $K,L,t\in\mathbb{N}$, let $R$ be a relation on natural numbers, and let
$\mathrm{Init}$ be a predicate on natural numbers. A natural number $u<K$ satisfies
$\mathrm{Init}(u)$ and is related to $t$ by a walk of length at most $L$ using $R$.
\end{definition}

\begin{definition}[Bounded reachability ball]
\lean{Complexity.LogProg.ISCount.ball}
\label{Complexity.LogProg.ISCount.ball}
\leanok
\uses{Complexity.LogProg.ISCount.Mem}
Let $K,L\in\mathbb{N}$, let $R$ be a relation on natural numbers, and let
$\mathrm{Init}$ and $\mathrm{Acc}$ be predicates on natural numbers. The ball at level
$L$ is the finite set of natural numbers $t<K$ for which there is some $u<K$ satisfying
$\mathrm{Init}(u)$ and connected to $t$ by a walk of length at most $L$ using $R$.
\end{definition}

\begin{lemma}[Bound along a bounded walk]
\lean{Complexity.LogProg.ISCount.walkLe_lt}
\label{Complexity.LogProg.ISCount.walkLe_lt}
\leanok
\uses{Complexity.WalkLe}
\difficulty{4}
Let $R$ be a binary relation on the natural numbers and let $K,j,c,t$ be natural
numbers. If $c<K$ and there is an $R$-walk from $c$ to $t$ of length at most $j$ whose
vertices are bounded by $K$, then $t<K$.
\end{lemma}

\begin{lemma}[Vertex bound for reachable states]
\lean{Complexity.LogProg.ISCount.Mem.lt}
\label{Complexity.LogProg.ISCount.Mem.lt}
\leanok
\uses{Complexity.LogProg.ISCount.Mem, Complexity.LogProg.ISCount.walkLe_lt}
\difficulty{1}
Let $n$ be a natural number, let $E$ be a relation on the natural numbers, and let
$K,L,t$ be natural numbers. Define membership by requiring that some initial vertex
$u<K$ satisfies the initial predicate and that $t$ is reachable from $u$ by a walk of
length at most $K$, with every step following $E$ through an intermediate vertex below
$n$. If $t$ belongs to this set for the specified $K,E,L$ and initial predicate, then
$t<K$.
\end{lemma}

\begin{lemma}[Membership in the reachable set]
\lean{Complexity.LogProg.ISCount.mem_ball}
\label{Complexity.LogProg.ISCount.mem_ball}
\leanok
\uses{Complexity.LogProg.ISCount.Mem, Complexity.LogProg.ISCount.Mem.lt, Complexity.LogProg.ISCount.ball}
\difficulty{3}
Let $K,L,t$ be natural numbers, let $R$ be a binary relation on the natural numbers, and
let $I$ be a predicate on the natural numbers. Define a walk of length at most $L$ from
$u$ to $t$ to allow stopping at the current vertex or taking a step along $R$, with each
vertex reached by a step required to be less than $K$. Then $t$ belongs to the set of
vertices reachable in at most $L$ steps from some $I$-vertex $u<K$ if and only if such a
starting vertex and walk exist.
\end{lemma}

\begin{lemma}[Monotonicity of reachable sets]
\lean{Complexity.LogProg.ISCount.ball_mono}
\label{Complexity.LogProg.ISCount.ball_mono}
\leanok
\uses{Complexity.LogProg.ISCount.ball, Complexity.LogProg.ISCount.mem_ball, Complexity.walkLe_mono}
\difficulty{3}
Let $K$ be a natural number, let $R$ be a binary relation on the natural numbers, and
let $\text{Init}$ be a predicate on the natural numbers. For each natural number $L$, let
the reachable set consist of those numbers below $K$ that can be reached in at most $L$
steps from some number below $K$ satisfying $\text{Init}$, where each step follows $R$ to
a number below $K$. If $L \le L'$, then the reachable set for $L$ is contained in the
reachable set for $L'$.
\end{lemma}

\begin{lemma}[One-step walk decomposition]
\lean{Complexity.LogProg.ISCount.walkLe_succ_iff}
\label{Complexity.LogProg.ISCount.walkLe_succ_iff}
\leanok
\uses{Complexity.WalkLe, Complexity.walkLe_refl, Complexity.walkLe_succ}
\difficulty{5}
Let $R$ be a binary relation on the natural numbers, and let $K,j,c,t$ be natural
numbers. There is an $R$-walk of length at most $j+1$ from $c$ to $t$, with intermediate
vertices less than $K$, if and only if either there is such a walk of length at most $j$
from $c$ to $t$, or there is a number $w$ such that there is such a walk of length at
most $j$ from $c$ to $w$, $R(w,t)$ holds, and $t<K$.
\end{lemma}

\begin{lemma}[Reachability after one additional step]
\lean{Complexity.LogProg.ISCount.mem_succ_iff}
\label{Complexity.LogProg.ISCount.mem_succ_iff}
\leanok
\uses{Complexity.LogProg.ISCount.Mem, Complexity.LogProg.ISCount.walkLe_succ_iff, Complexity.walkLe_succ}
\difficulty{4}
Let $K,L,t$ be natural numbers, let $R$ be a binary relation on the natural numbers, and
let $I$ be a predicate on the natural numbers. A vertex $t$ is reachable in at most
$L+1$ steps from some $I$-vertex below $K$ if and only if it is reachable in at most $L$
steps, or there is a vertex $w$ reachable in at most $L$ steps such that $R(w,t)$ and
$t<K$.
\end{lemma}

\begin{lemma}[Nonmembership after one additional step]
\lean{Complexity.LogProg.ISCount.not_mem_succ_iff}
\label{Complexity.LogProg.ISCount.not_mem_succ_iff}
\leanok
\uses{Complexity.LogProg.ISCount.Mem, Complexity.LogProg.ISCount.ball, Complexity.LogProg.ISCount.mem_ball, Complexity.LogProg.ISCount.mem_succ_iff}
\difficulty{4}
Let $K,L,v$ be natural numbers, let $R$ be a binary relation on the natural numbers, and
let $I$ be a predicate on the natural numbers, and suppose $v<K$. The vertex $v$ is not
reachable in at most $L+1$ steps from any $I$-vertex below $K$ if and only if every
vertex $u$ reachable in at most $L$ steps from such a starting vertex is distinct from
$v$ and does not satisfy $R(u,v)$.
\end{lemma}

\begin{lemma}[Initial reachable set has size one]
\lean{Complexity.LogProg.ISCount.card_ball_zero}
\label{Complexity.LogProg.ISCount.card_ball_zero}
\leanok
\uses{Complexity.LogProg.ISCount.ball, Complexity.LogProg.ISCount.mem_ball}
\difficulty{4}
Let $K$ be a natural number, let $R$ be a binary relation on the natural numbers, and
let $\text{Init}$ be a predicate on the natural numbers. If $u_0<K$ and $\text{Init}(u)$
holds exactly when $u=u_0$, then the set of numbers below $K$ reachable in zero steps
from an initial number has cardinality $1$.
\end{lemma}

\begin{definition}[Recursive branching process]
\lean{Complexity.LogProg.ISCount.GS}
\label{Complexity.LogProg.ISCount.GS}
\leanok
Let $A,B,\mathrm{Kont}$ be predicates on natural numbers. For zero rounds, the process
at position $t$ and count $c$ holds exactly when $\mathrm{Kont}(c)$ holds. For $r+1$
rounds, it holds when either $A(t)$ holds and the process continues for $r$ rounds at
position $t+1$ with count $c+1$, or $B(t)$ holds and it continues at position $t+1$ with
count $c$.
\end{definition}

\begin{lemma}[Conjunction at the terminal condition]
\lean{Complexity.LogProg.ISCount.gs_and}
\label{Complexity.LogProg.ISCount.gs_and}
\leanok
\uses{Complexity.LogProg.ISCount.GS}
\difficulty{3}
Let $A$, $B$, and $\text{Kont}$ be predicates on the natural numbers, let $Q$ be a
proposition, and let $r,t,c$ be natural numbers. The recursive condition that makes $r$
choices between $A$ and $B$, ending with the condition $\text{Kont}(c')\land Q$, holds at
$(r,t,c)$ if and only if the same recursive condition ending with $\text{Kont}(c')$ holds
there and $Q$ holds.
\end{lemma}

\begin{lemma}[Finite-set characterization of the recursive condition]
\lean{Complexity.LogProg.ISCount.gs_iff}
\label{Complexity.LogProg.ISCount.gs_iff}
\leanok
\uses{Complexity.LogProg.ISCount.GS}
\difficulty{6}
Let $A$, $B$, and $\text{Kont}$ be predicates on the natural numbers, and let $r,t,c$ be
natural numbers. The recursive condition that makes $r$ successive choices between $A$
and $B$, incrementing the counter exactly when $A$ is chosen and ending with
$\text{Kont}$, holds at $(r,t,c)$ if and only if there is a finite set
$T\subseteq\{t,t+1,\ldots,t+r-1\}$ such that $A(y)$ holds for every $y\in T$, $B(y)$
holds for every other $y$ in that interval, and $\text{Kont}(c+|T|)$ holds.
\end{lemma}

\begin{lemma}[Inner-stage condition on the reachable set]
\lean{Complexity.LogProg.ISCount.gs_inner}
\label{Complexity.LogProg.ISCount.gs_inner}
\leanok
\uses{Complexity.LogProg.ISCount.GS, Complexity.LogProg.ISCount.Mem, Complexity.LogProg.ISCount.Mem.lt, Complexity.LogProg.ISCount.ball, Complexity.LogProg.ISCount.gs_iff, Complexity.LogProg.ISCount.mem_ball}
\difficulty{4}
Let $K$ be a natural number, let $R$ be a binary relation on the natural numbers, and
let $\text{Init}$ be a predicate on the natural numbers. For natural numbers $L$ and $v$,
the recursive condition over $K$ stages, starting at position $0$ with counter $0$,
holds—requiring at each selected position $u$ that $u$ is reachable within $L$ steps
from an initial number, that $u\ne v$, and that $R(u,v)$ fails; allowing the alternative
condition to hold everywhere; and ending when the counter equals the cardinality of the
reachable set—if and only if every element $u$ of that reachable set satisfies $u\ne v$
and not $R(u,v)$.
\end{lemma}

\begin{lemma}[Outer-stage condition on the reachable set]
\lean{Complexity.LogProg.ISCount.gs_outer}
\label{Complexity.LogProg.ISCount.gs_outer}
\leanok
\uses{Complexity.LogProg.ISCount.GS, Complexity.LogProg.ISCount.Mem, Complexity.LogProg.ISCount.Mem.lt, Complexity.LogProg.ISCount.ball, Complexity.LogProg.ISCount.gs_iff, Complexity.LogProg.ISCount.mem_ball}
\difficulty{5}
Let $K$ and $L$ be natural numbers, let $R$ be a binary relation on the natural numbers,
and let $\text{Init}$ and $\text{Acc}$ be predicates on the natural numbers. Suppose
$B(t)$ holds exactly when $t$ is not reachable within $L$ steps from an initial number,
for every $t<K$. The recursive condition over $K$ stages, starting at position $0$ with
counter $0$, holds—requiring at each selected position that it be reachable and fail
$\text{Acc}$, using $B$ at the alternative positions, and ending with $\text{Kont}$—if and
only if every element of the reachable set fails $\text{Acc}$ and $\text{Kont}$ holds at
the cardinality of that set.
\end{lemma}
